% Options for packages loaded elsewhere
\PassOptionsToPackage{unicode}{hyperref}
\PassOptionsToPackage{hyphens}{url}
\documentclass[
  ukrainian,
]{article}
\usepackage{xcolor}
\usepackage[margin=2.5cm]{geometry}
\usepackage{amsmath,amssymb}
\setcounter{secnumdepth}{-\maxdimen} % remove section numbering
\usepackage{iftex}
\ifPDFTeX
  \usepackage[T1]{fontenc}
  \usepackage[utf8]{inputenc}
  \usepackage{textcomp} % provide euro and other symbols
\else % if luatex or xetex
  \usepackage{unicode-math} % this also loads fontspec
  \defaultfontfeatures{Scale=MatchLowercase}
  \defaultfontfeatures[\rmfamily]{Ligatures=TeX,Scale=1}
\fi
\usepackage{lmodern}
\ifPDFTeX\else
  % xetex/luatex font selection
\fi
% Use upquote if available, for straight quotes in verbatim environments
\IfFileExists{upquote.sty}{\usepackage{upquote}}{}
\IfFileExists{microtype.sty}{% use microtype if available
  \usepackage[]{microtype}
  \UseMicrotypeSet[protrusion]{basicmath} % disable protrusion for tt fonts
}{}
\makeatletter
\@ifundefined{KOMAClassName}{% if non-KOMA class
  \IfFileExists{parskip.sty}{%
    \usepackage{parskip}
  }{% else
    \setlength{\parindent}{0pt}
    \setlength{\parskip}{6pt plus 2pt minus 1pt}}
}{% if KOMA class
  \KOMAoptions{parskip=half}}
\makeatother
\usepackage{graphicx}
\graphicspath{{/Users/kalinicenkovitalijmikolajovic/.codex/worktrees/5c15/Script_for_translate_book/Output_Final/farlou_chunked_v1_pdf_chunks/}}
\makeatletter
\newsavebox\pandoc@box
\newcommand*\pandocbounded[1]{% scales image to fit in text height/width
  \sbox\pandoc@box{#1}%
  \Gscale@div\@tempa{\textheight}{\dimexpr\ht\pandoc@box+\dp\pandoc@box\relax}%
  \Gscale@div\@tempb{\linewidth}{\wd\pandoc@box}%
  \ifdim\@tempb\p@<\@tempa\p@\let\@tempa\@tempb\fi% select the smaller of both
  \ifdim\@tempa\p@<\p@\scalebox{\@tempa}{\usebox\pandoc@box}%
  \else\usebox{\pandoc@box}%
  \fi%
}
% Set default figure placement to htbp
\def\fps@figure{htbp}
\makeatother
\ifLuaTeX
\usepackage[bidi=basic,shorthands=off,provide=*]{babel}
\else
\usepackage[bidi=default,shorthands=off,provide=*]{babel}
\fi
\ifLuaTeX
  \usepackage{selnolig} % disable illegal ligatures
\fi
\setlength{\emergencystretch}{3em} % prevent overfull lines
\providecommand{\tightlist}{%
  \setlength{\itemsep}{0pt}\setlength{\parskip}{0pt}}
\usepackage{bookmark}
\IfFileExists{xurl.sty}{\usepackage{xurl}}{} % add URL line breaks if available
\urlstyle{same}
\hypersetup{
  pdflang={uk},
  hidelinks,
  pdfcreator={LaTeX via pandoc}}

\author{}
\date{}

\begin{document}

\section{Обчислення
інтеграла}\label{ux43eux431ux447ux438ux441ux43bux435ux43dux43dux44f-ux456ux43dux442ux435ux433ux440ux430ux43bux430}

Припустимо, ми хочемо обчислити інтеграл

Здесь \emph{k} --- модуль Юнга, характеризующий упругость материала.
Материалы, обладающие большим модулем Юнга, колеблются на более высоких
частотах. Звуковые волны представляют собой продольные волны.

\begin{verbatim}
2. Если линейная плотность струны   \rho(x)   зависит от координаты, то волновое урагнение записывается в виде
\end{verbatim}

(детерміноване завдання). Щоб використати метод Монте-Карло, потрібно
винайти гру на удачу, яка дає приблизне значення цього інтеграла.

Звісно, можна придумати багато різних ігор. Остаточна версія залежить
від точності розрахунків, простоти гри тощо. Очевидний варіант гри для
обчислення інтегралу --- кинути дротик у прямокутний
\(R = \{(x,y): a \le x \le b, 0 \le y \le \max f(x)\}\) (див. рис. 42.2)

\pandocbounded{\includegraphics[keepaspectratio]{images/_page_315_Figure_7.jpeg}}

Рис. 42 2. Обчислення інтегралу методом Монте-Карло.

Очевидно, якщо випадковим чином кинути близько 100 дротиків у
прямокутник R, що містить графік функції, можна отримати оцінку значення
інтеграла, помноживши площу прямокутника R на відносну кількість
дротиків під графіком функції. Отже, результат нашої гри

т. е. получается УЧП с переменными коэффициентами.

\begin{verbatim}
3. Поскольку волновое уравнение   u_{tt} = \alpha^2 u_{xx}   содержит производную по времени второго порядка, для получения единственного решения при t>0 необходимо задать   \partial sa   начальных условия:
\end{verbatim}

--- це інтегральний бал 1.

Щоб виконати ці обчислення на комп'ютері, нам потрібно якимось чином
отримати послідовність випадкових точок (подивимось, як це буде зроблено
пізніше), тобто наказати комп'ютеру «кидати дротики». Блок-схема на рис.
42.3 показує, як комп'ютер вирішить цю проблему.

\end{document}
